\documentclass[a4paper,12pt]{article}

\usepackage[utf8]{inputenc}
\usepackage[T1]{fontenc}
\usepackage[polish]{babel}
\usepackage{amsmath, amssymb}
\usepackage{graphicx}
\usepackage{float}
\usepackage{hyperref}
\usepackage{geometry}
\usepackage{csvsimple} 

\geometry{margin=2.5cm}

\title{Sprawozdanie 1 z Statystyki i Modeli Liniowych}
\author{Wiktor Zoga}
\date{\today}

\begin{document}

\maketitle
\newpage

\section{Zadanie 1 - Porównywanie estymatorów}

\subsection{Opis eksperymentu}
W tym zadaniu wygenerowano $n$ obserwacji z rozkładu normalnego $N(\theta, \sigma^2)$ dla następujących zestawów parametrów:

\begin{enumerate}
    \item $\theta = 0$, $\sigma = 1$
    \item $\theta = 0$, $\sigma = 2$
    \item $\theta = 4$, $\sigma = 1$
    \item $\theta = 4$, $\sigma = 2$
\end{enumerate}

Dla każdej próbki obliczono siedem estymatorów wartości oczekiwanej:

\begin{itemize}
    \item $\hat{\theta}_1$ -- średnia arytmetyczna
    \item $\hat{\theta}_2$ -- mediana
    \item $\hat{\theta}_3$ -- ważona średnia z własnymi wagami
    \item $\hat{\theta}_4$ -- estymator z wagami zależnymi od dystrybuanty normalnej
    \item $\hat{\theta}_5$ -- średnia harmoniczna
    \item $\hat{\theta}_6$ -- średnia potęgowa rzędu 3
    \item $\hat{\theta}_7$ -- własny estymator (średnia z minimum i maksimum próby)
\end{itemize}

Doświadczenie powtórzono $10{,}000$ razy dla każdego zestawu parametrów i rozmiarów próbek $n = 20, 50, 100$.

\newpage

\section{Wyniki - wykresy pudełkowe i tabele statystyk}

\newcommand{\offplotwithstats}[3]{
    \begin{figure}[H]
        \centering
        \includegraphics[width=0.9\textwidth]{../plots/5_off_boxplot_theta#1_sigma#2_n#3.png}
        \caption{Boxplot estymatorów bez piątego.}
    \end{figure}
}

% --- makro do wstawienia wykresu i tabeli CSV pod nim
\newcommand{\plotwithstats}[3]{
    \subsection{$\theta=#1$, $\sigma=#2$, $n=#3$}
    \begin{figure}[H]
        \centering
        \includegraphics[width=0.9\textwidth]{../plots/boxplot_theta#1_sigma#2_n#3.png}
        \caption{Boxplot estymatorów dla $\theta=#1$, $\sigma=#2$, $n=#3$.}
    \end{figure}

    \offplotwithstats{#1}{#2}{#3}

    \begin{center}
    \csvreader[
        tabular=|l|c|c|c|c|,
        table head=\hline Estymator & Średnia & Wariancja & Bias & MSE \\ \hline,
        late after line=\\\hline,
        filter expr={true},  % brak filtrowania
        respect percent
    ]{../plots/stats_theta#1_sigma#2_n#3.csv}{Estimator=\est,Mean=\mean,Variance=\var,Bias=\bias,MSE=\mse}
    {\est & \pgfmathprintnumber{\mean} & \pgfmathprintnumber{\var} & \pgfmathprintnumber{\bias} & \pgfmathprintnumber{\mse} }
    \end{center}

    % \newpage
}

% --- teraz wywołanie dla wszystkich kombinacji parametrów

\plotwithstats{0}{1}{20}
\plotwithstats{0}{1}{50}
\plotwithstats{0}{1}{100}

\plotwithstats{0}{2}{20}
\plotwithstats{0}{2}{50}
\plotwithstats{0}{2}{100}

\plotwithstats{4}{1}{20}
\plotwithstats{4}{1}{50}
\plotwithstats{4}{1}{100}

\plotwithstats{4}{2}{20}
\plotwithstats{4}{2}{50}
\plotwithstats{4}{2}{100}

\subsection{Dyskusja}

\subsubsection{Est1 - Średnia Arytmetyczna}
Zgodnie z teorią estymacji, średnia arytmetyczna jest nieobciążonym estymatorem średniej w rozkładzie normalnym.
Wyniki symulacji to potwierdzają – $\hat{\theta}_1$ wykazuje obciążenie bliskie zeru 
oraz najniższą wariancję (i tym samym najniższe MSE) spośród wszystkich badanych estymatorów nieobciążonych. 
Służy jako dobry punkt odniesienia.

\subsubsection{Est2 - Mediana}
Mediana również jest estymatorem nieobciążonym dla rozkładu normalnego. 
Wynika to z tego, że rozkład ma symetryczną gęstość, a mediana ma "własność przesunięcia".
Jej wariancja jest jednak wyższa niż wariancja średniej arytmetycznej,
co potwierdzają wyniki MSE. 
Dla średniej rozkładu normalnego, jest to estymator mniej efektywny niż średnia arytmetyczna, 
choć bardziej odporny na wartości odstające (co nie było przedmiotem tego badania).

\subsubsection{Est3 - Ważona Średnia}
Ważona średnia (losowe wagi o sumie jeden) jest estymatorem nieobciążonym, który może być bardziej efektywny niż średnia arytmetyczna w przypadku rozkładów o asymetrycznych ogonach.
W przypadkach rozkładu normalnego, nie gra to większej roli, a MSE jest lekko wyższy niż dla średniej arytmetycznej.

\subsubsection{Est4 - Ważona Średnia wagami z Dystrybuanty Normalnej}
Ten estymator jest obciążony ponieważ wagi nie sumują się do 1. Po mimo zauważalnie mniejszej wariancji 
od innych estymatorów ma duży bias co przekłada się na wysokie MSE.

\subsubsection{Est5 - Średnia Harmoniczna}
Średnia harmoniczna jest estymatorem obciążonym, szczególnie w przypadku rozkładów z dużą wariancją.
Wyniki pokazują znaczne obciążenie i wysokie MSE, co czyni ją nieodpowiednią do estymacji średniej w rozkładzie normalnym.
Wynika to głównie z faktu, że $\frac{1}{X}$ nie jest funkcją całkowaną z gestością rozkładu normalnego.

\subsubsection{Est6 - Średnia Potęgowa Rzędu 3}
Dzięki mocnemu prawu wielkich liczb średnia potęgowa rzędu 3 jest estymatorem asymptotycznie nieobciążonym, ale jedynie dla średniej 0.
W przypadkach z $\theta = 4$ jest zauważalny bias.
Wariancja jest jednak wyższa niż w przypadku średniej arytmetycznej, co czyni ją mniej efektywną w praktyce.

\subsubsection{Est7 - Średnia z Minimum i Maksimum Próby}
Ten estymator jest nieobciążony, co wynika z bezpośredniej symetrii rozkładu normalnego (pomijam rachunki, ale wykonałem je).
Jednakże jego wariancja jest znacznie wyższa niż w przypadku innych estymatorów, gdyż nie maleje wraz ze wzrostem $n$ jak w przypadku innych estymatorów.
Teoretyczne wyliczenia wariancji nie są już tak proste, jednak wyniki symulacji pokazują, że jest ona wyższa niż w przypadku średniej, ale wciąż nie duża.

\end{document}
